\section{The principal propositions}\label{subsec:principal}

The five propositions below use the joint choices described in Section~\ref{sec:strategy}.
Their complete proofs belong to Part~\ref{part:construction}; the explanations here give the ideas needed to follow the deduction of Theorem~\ref{thm:main}.

\paragraph{Removing the lattice spacing.}
At a fixed physical volume, increasing the number of sites must not change what the model looks like at a fixed larger distance.
The coupling and the normalisation of each observable are chosen to achieve this compatibility.
The first proposition gives both one such choice and convergence along its entire spacing sequence.

\begin{proposition}[Ultraviolet construction]\label{prop:uv}
For every gauge group with simple $\g$ (Definition~\ref{def:gauge}), there exists a regular Wilson family $\W$ (Definitions~\ref{def:family} and~\ref{def:regular}) such that, for every $p,n\in\N$ and $v\in\mathcal T^n$, there exists a smooth kernel $K_{p,v}$ on the periodically separated set of Definition~\ref{def:periodic} with
\begin{equation}\label{eq:fixed-volume}
\lim_{N\to\infty}C_{N,p}(v;f)=\int f(x)K_{p,v}(x)\,dx
\end{equation}
for every separated test $f$ supported inside volume $p$ (Definitions~\ref{def:moments}, \ref{def:tests} and~\ref{def:admissible}).
\end{proposition}

\emph{Idea of the proof.}
Replace the fine link variables by averaged variables at a larger spacing, keeping the integral of the probability density exactly the same.
Where the fields are small, their fluctuations can be compared with a Gaussian integral, whose energy is quadratic.
Where they are large, the energy suppresses their contribution; these configurations are estimated rather than discarded.
Repeating this integration changes the coefficient multiplying the squared curvature.
We account for that change in the coupling, instead of treating it as a small error.
The resulting coupling grows logarithmically as the original spacing shrinks.

We carry the local observables through the same integrations.
Their normalisation is fixed at a reference scale, so different cutoffs are compared there rather than at their different bare couplings.
To prove convergence, stop the integration at an intermediate scale.
At this fixed scale the effective correlations converge with the cutoff, while the remaining error is uniformly small when the intermediate scale is sufficiently fine.
Choosing the intermediate scale first and the cutoffs second proves the Cauchy property.

\gap{The later ultraviolet sections must prove the estimates for integrating small and large fields, identify the change of the actual effective coupling, and establish the fixed-scale Cauchy estimate with a uniformly small remaining error.
The fixed-volume conclusion has been checked in the formal source; its analytic proof is still missing from this manuscript.\sourcefootnote{\nolinkurl{Construction.lean}, lines 15--40; \nolinkurl{HierarchyBridge.lean}, lines 16--26; \nolinkurl{YM/FullVolume/Profiles.lean}, lines 25--47.}}

\paragraph{Enlarging the volume and separating phases.}
The estimates for a fixed collection of observables must remain bounded when the surrounding box grows.
They give compactness, but the limit can retain several possible vacuum states.
The next proposition includes the separate argument that resolves this average into phases.

\begin{proposition}[Continuum limit and phases]\label{prop:limit}
For the jointly chosen family of Proposition~\ref{prop:uv} and every strictly increasing $p:\N\to\N$, there exist a strictly increasing $q:\N\to\N$ and hierarchy data (Definition~\ref{def:hierarchy}) satisfying the limit~\ref{main:1} and the phase and regularity conditions~\ref{recon:7}--\ref{recon:12} of Definition~\ref{def:reconstructed}.
\end{proposition}

\emph{Idea of the proof.}
Reflection positivity bounds an average of several separated observables by averages of reflected copies of each individual observable.
The ultraviolet estimates bound these individual averages independently of box size.
Integrating between successive lattice slices controls changes in the correlations when their points move, as explained in Subsection~\ref{subsec:iv-compactness}.
Uniform bounds on the values and on these changes allow a subsequence to be selected for all polynomial basis tuples and all compact sets at once (Lemma~\ref{lem:uv-labels}).
Lattice summation by parts relates derivatives of the limits to insertions of higher degree and gives smoothness.
The fixed-volume limits of Proposition~\ref{prop:uv} identify every inner limit, so only the volume subsequence remains to be chosen.

The average state is split into phases by examining observables separated by a very long time.
In the Hilbert space obtained from reflection positivity, long-time evolution removes the positive-energy part and retains the zero-energy vectors.
Observables placed between these long evolutions act on that zero-energy space.
Once rotations are established, exchanging a time direction with a spatial direction shows that these limiting operators commute.
Representing these commuting operators by continuous functions gives the compact phase space; their vacuum expectations give its probability measure.
This produces the mixture identity~\eqref{eq:mixture}.

\gap{The infinite-volume argument still needs the same-family bounds for local observables and their derivative representatives from the ultraviolet sections.
The later symmetry and phase sections must also prove that the contribution from a far-away stress integral tends to zero and that the long-time operators commute; neither follows from compactness alone.\sourcefootnote{\nolinkurl{YM/Extraction/Limits.lean}, lines 26--43; \nolinkurl{HierarchyBridge.lean}, lines 44--57 and 82--119; the phase representation is in \nolinkurl{YM/Phase/Mixture.lean}.}}

\paragraph{Recovering relativistic fields.}
Euclidean time evolution damps positive energies.
Allowing time to become complex relates this damping to unitary time evolution, and hence to relativistic vacuum expectations.
The estimates must hold as the imaginary time separations tend to zero, because the resulting fields are distributions.

\begin{proposition}[Relativistic reconstruction]\label{prop:reconstruction}
For the jointly chosen hierarchy of Proposition~\ref{prop:limit}, the selected theory satisfies~\ref{recon:1}--\ref{recon:6} and the continuation satisfies~\ref{recon:13}--\ref{recon:15} of Definition~\ref{def:reconstructed}, with the same fields and phase as in Definition~\ref{def:hierarchy}.
\end{proposition}

\emph{Idea of the proof.}
First place Euclidean observables in thin, separated time intervals.
Their bounds show that the continued correlations have distributional limits when the imaginary time separations disappear.
Positivity turns these limits into inner products of field words applied to the vacuum, with words as in~\eqref{eq:domain}.
Completing this space gives the Hilbert space, and multiplication by one more smeared field defines its operator action.

Rotation identities continue to Lorentz transformations.
Together with positivity of Euclidean energy, this confines energy and momentum to the forward cone.
To prove locality, compare the two orders of neighbouring fields at spacelike separated points: their analytic continuations agree on the Euclidean configurations and therefore give the same boundary values there.
The quotient in Proposition~\ref{prop:iv-smooth} removes classical polynomials whose correlations all vanish; the phase construction retains these null relations.
This gives faithful field labels in the sense of~\ref{recon:2}, while the constant observable $1$ supplies a nonzero field.

\gap{The reconstruction section must prove the bounds as imaginary time separations vanish and justify the equality of the two analytic orderings at spacelike separation.
The exact selected-phase and product-test conclusions have been checked in the formal source; these analytic arguments remain unwritten.\sourcefootnote{\nolinkurl{PublicTheory.lean}, lines 93--319; \nolinkurl{HierarchyBridge.lean}, lines 82--133; boundary and locality proofs are organised in \nolinkurl{YM/Work/C01.lean} through \nolinkurl{C29.lean}.}}

\paragraph{Replacing a cluster by a single field.}
A distant observation cannot resolve every detail of several nearby insertions.
The operator product expansion expresses their effect through a finite list of fields at the cluster centre.
The coefficients must be independent of which distant observations are subsequently made.

\begin{proposition}[Common operator product expansion]\label{prop:ope}
The fields and coefficients in~\ref{hierarchy:7} of the jointly chosen hierarchy satisfy~\ref{ope:1}--\ref{ope:2} of Definition~\ref{def:ope}, for the clusters of Definition~\ref{def:cluster}.
\end{proposition}

\emph{Idea of the proof.}
Average the clustered observables from their separation scale toward a larger scale, keeping track of which single local fields they produce.
Retain the terms through degree $D$ and estimate everything else.
The first omitted degree gives the power $D+1$ in the remainder.
The arbitrarily small loss in that power accommodates logarithmic factors.
A fixed collection of correlations with distant test fields determines the coefficients once, before any later spectator fields or phase are chosen.
The error has an operator-norm bound before these commuting operators are evaluated in a phase, so the same bound holds in every phase and in the average state.

\gap{The product section must prove the local-product remainder estimate and show how the fixed coefficient choice gives an operator bound before evaluation in a phase.
The quantifiers and both powers of the radius in~\ref{ope:1}--\ref{ope:2} have been checked; their analytic proof is still deferred.\sourcefootnote{\nolinkurl{ChallengeSpec.lean}, lines 58--80; \nolinkurl{HierarchyBridge.lean}, lines 68--80; \nolinkurl{YM/OPE/Orbits/Kernel.lean}, lines 62--77.}}

\paragraph{Detecting the logarithmic correction.}
The construction supplies a stress tensor and subtracts the selected vacuum expectation from its trace.
Comparison at short distances identifies the leading two-point singularity of this centred field.
The logarithmic correction to that singularity is incompatible with the finite free-field expressions in Definition~\ref{def:wick}.

\begin{proposition}[Trace law and finite-profile exclusion]\label{prop:trace}
The field $\Theta\in F_4$ of the jointly chosen hierarchy satisfies~\ref{trace:1}--\ref{trace:4} of Definition~\ref{def:trace} in the theory of Proposition~\ref{prop:reconstruction}, with the two-point function of Definition~\ref{def:trace-pair}.
\end{proposition}

\emph{Idea of the proof.}
At small distances the fields can be compared, after adjusting their normalisations, with the free Maxwell fields of Section~\ref{sec:maxwell}.
The trace carries an additional inverse running coupling at each insertion.
Since that coupling grows logarithmically, its two-point function acquires the inverse square logarithm in~\eqref{eq:trace-law}.
The coefficient of the logarithmic growth enters both the running coupling and the trace normalisation, and cancels from their ratio.
The final amplitude by itself therefore does not check that coefficient.

The comparison initially works at specially chosen scales.
The positive-energy spectral representation expresses the two-point function through decaying exponentials of the time separation with positive weights.
It therefore decreases with separation, so its values just above and below a requested separation bound its value there.
Every sequence of radii approaching zero has a subsequence where the comparison applies; thus a sequence contradicting the full limit cannot exist.
Finite Wick profiles instead have expansions in integer powers of $r$ and nonnegative powers of $\log r$.
Such expansions cannot have the positive inverse-logarithmic behaviour in~\eqref{eq:trace-law}.

\gap{The trace section must prove the stress normalisation, comparison along subsequences of arbitrary radii, the spectral comparison between nearby separations, and the small-radius expansions of the finite comparison kernels~\eqref{eq:mode}.
The amplitude and the exact finite comparison class have been checked; these proofs remain to be supplied.\sourcefootnote{\nolinkurl{TraceBridge.lean}, lines 14--44; \nolinkurl{YM/TraceExtension/Verse/Cape.lean}, lines 53--88; \nolinkurl{YM/TraceExtension/Wick/Log.lean}, lines 96--108.}}

\begin{proof}[Deduction of Theorem~\ref{thm:main} from the principal propositions]
Fix the gauge group and choose the family of Proposition~\ref{prop:uv}, including its insertions and positive regularity slope.
For a prescribed strictly increasing $p$, take the common $q$ and hierarchy of Proposition~\ref{prop:limit}.
This gives~\ref{main:1} and~\ref{recon:7}--\ref{recon:12} on that refinement.
Proposition~\ref{prop:reconstruction} supplies the remaining reconstruction properties, hence~\ref{main:2}.
Propositions~\ref{prop:ope} and~\ref{prop:trace} give~\ref{main:3} and~\ref{main:4} for those same objects.
\end{proof}

